En esta sección se presentan dos soluciones al problema de encontrar las diferencias mínimas entre dos secuencias, siguiendo dos paradigmas de diseño de algoritmos vistos en el curso: fuerza bruta y programación dinámica. Ambas soluciones buscan determinar las posiciones donde las cadenas difieren, basándose en la subsecuencia común más larga (LCS) como estructura intermedia.

A continuación, se describe cada algoritmo con su pseudocódigo respectivo, junto a un análisis de complejidad temporal y espacial en función de las longitudes de entrada $n = |s|$ y $m = |t|$.
